\chapter{Topological Sigma Models}
\label{ch:tqft-topological-sigma-models}

\ConventionsEuclidean

Let \(\Sigma\) be a compact oriented smooth surface equipped with a complex
structure \(j_\Sigma\), and let \(X\) be a compact smooth manifold.  A
two-dimensional sigma model has bosonic field \(\phi:\Sigma\to X\).  A
topological sigma model is a cohomological sigma model in the sense of
Chapter~\ref{ch:tqft-cohomological-field-theories}: it has an odd differential
\(Q\), \(Q\)-closed observables, and a stress-tensor or metric-response
insertion which is \(Q\)-exact after the anomaly and boundary terms have been
accounted for.  The central examples are the A-model and the B-model obtained
from two-dimensional \(\mathcal N=(2,2)\) sigma models by two distinct twists.

This chapter develops the finite-dimensional geometry which the path integral
is expected to produce after localization.  The identification of a
Gromov--Witten integral with a continuum QFT is made at the level where the
regulator, integration cycle, compactification, anomaly line, and gluing maps
have been specified.  Those data are part of the QFT problem surrounding the
enumerative formula.

\begin{hypothesis}[A-model cohomological sigma-model datum]
\label{hyp:a-model-cohomological-sigma-datum}
A closed A-model datum is a tuple
\[
  \mathfrak A_X
  =
  \bigl(
    \Sigma,j_\Sigma; X,\omega,J,g_X; [B+\ii\omega];
    \mathcal F_{\Sigma,X},Q_A,\bar\partial_J,
    \mathfrak V_X,\Lambda_{\rm Nov},\mathfrak c_A
  \bigr)
\]
with the following entries.
\begin{enumerate}[label=\textnormal{(\roman*)}]
\item \((\Sigma,j_\Sigma)\) is a compact marked Riemann surface, and
  \((X,\omega,J,g_X)\) is compact almost K\"ahler.
\item \(\mathcal F_{\Sigma,X}\) is the locally super-ringed mapping field
  space whose bosonic body is \(C^\infty(\Sigma,X)\) and whose odd
  directions are the twisted fermionic section spaces.
\item \(Q_A\) is an odd cohomological vector field whose scalar local
  observable complex is identified with \((\Omega^\bullet(X),d)\), and
  \(\bar\partial_J\phi=0\) is the localization equation for the bosonic
  zero modes.
\item The action splits into a \(Q_A\)-exact nonnegative term plus the
  complexified K\"ahler weight
  \(Q^\beta=\exp(2\pi\ii\int_\beta(B+\ii\omega))\) on each curve class
  \(\beta\in H_2(X;\mathbb Z)\).
\item \(\mathfrak V_X\) is a virtual integration package for the stable-map
  compactifications \(\overline{\mathcal M}_{g,n}(X,\beta)\), including
  evaluation maps, forgetful maps, deformation identifications, and gluing
  compatibility.
\item \(\Lambda_{\rm Nov}\) is the Novikov completion in which the instanton
  sum is interpreted, and \(\mathfrak c_A\) is the source, collision, and
  contact-term prescription used to define products of local insertions.
\end{enumerate}
The datum is a hypothesis about the cohomological QFT layer.  The
finite-dimensional Gromov--Witten package \(\mathfrak V_X\) is an output
expected from localization once this datum is constructed.  The continuum QFT
construction itself is the separate problem recorded in
Open Problem~\ref{op:topological-sigma-model-continuum-qft}.
\end{hypothesis}

\begin{hypothesis}[Closed B-model cohomological sigma-model datum]
\label{hyp:b-model-cohomological-sigma-datum}
A closed B-model datum is a tuple
\[
  \mathfrak B_X
  =
  \bigl(
    \Sigma,j_\Sigma; X,J_X,\Omega_X;
    \mathcal F_{\Sigma,X},Q_B,{\rm PV}^{\bullet,\bullet}(X),
    {\rm Tr}_{\Omega_X},\mathfrak c_B
  \bigr)
\]
where \(X\) is a compact complex \(m\)-fold with trivial canonical bundle,
\(\Omega_X\) is a nowhere-vanishing holomorphic volume form,
\({\rm PV}^{p,q}(X)=\Omega^{0,q}(X,\wedge^pT^{1,0}X)\), \(Q_B=\bar\partial\)
on scalar local observables, and
\[
  {\rm Tr}_{\Omega_X}(\mu)
  =
  \int_X \Omega_X\wedge\iota_\mu\Omega_X
\]
is the trace pairing when the total polyvector and Dolbeault degrees make the
integrand top degree.  The datum also includes the anomaly-free axial twist,
the regulated \(Q_B\)-invariant integration cycle, and the contact-term
prescription \(\mathfrak c_B\).  Its first-order deformation complex is the
Dolbeault polyvector differential graded Lie algebra
\[
  \bigl({\rm PV}^{1,\bullet}(X),\bar\partial,[\cdot,\cdot]\bigr),
\]
with Maurer--Cartan equation
\(\bar\partial\mu+\frac12[\mu,\mu]=0\).
\end{hypothesis}

\section{The Mapping Field Space}

\begin{definition}[Mapping space and fermionic directions]
Let \(X\) be a smooth manifold.  The smooth bosonic field space is the
Fr\'echet manifold
\[
  \mathcal F_{\Sigma,X}^{\rm bos}
  =
  C^\infty(\Sigma,X).
\]
At \(\phi\in C^\infty(\Sigma,X)\), its tangent space is
\[
  T_\phi \mathcal F_{\Sigma,X}^{\rm bos}
  =
  \Gamma(\Sigma,\phi^\ast TX).
\]
A fermionic field in a sigma model is a parity-reversed section of a vector
bundle built functorially from \(T\Sigma\), \(T^\ast\Sigma\), and \(\phi^\ast TX\).
Thus the full field space is a locally super-ringed Fr\'echet space over
\(C^\infty(\Sigma,X)\).  Its structure sheaf is locally the completed exterior
algebra on the continuous dual of the chosen fermionic section spaces.
\end{definition}

The last sentence is part of the definition.  Fermionic field variables are
odd generators in a supercommutative algebra of functions on field space,
rather than elements of the underlying bosonic mapping space.  Fermionic field
operators in Lorentzian QFT are linear operators on a Hilbert-space domain and
are different objects.

For the A-model datum of
Hypothesis~\ref{hyp:a-model-cohomological-sigma-datum}, the target is an
almost K\"ahler manifold
\((X,\omega,J,g_X)\), where \(J^2=-1\),
\[
  \omega(u,v)=g_X(Ju,v),
  \qquad
  d\omega=0,
\]
and \(J\) is compatible with \(\omega\).  For the B-model datum of
Hypothesis~\ref{hyp:b-model-cohomological-sigma-datum}, the target is a
compact complex manifold \(X\) with trivial canonical bundle and a choice of
nowhere-vanishing holomorphic volume form \(\Omega_X\); a K\"ahler metric is
used to write a Lagrangian representative, but the cohomological observables
depend on complex-structure data.

\section{The A-Model Complex}

The scalar A-model differential acts on local functions of \(\phi\) as the de
Rham differential on \(X\).  In local coordinates \(x^i\) on \(X\), introduce
an odd scalar field \(\psi^i\) transforming as \(d x^i\).  On zero-form
observables,
\[
  Q_A \phi^i=\psi^i,
  \qquad
  Q_A\psi^i=0 .
\]
For \(\alpha\in\Omega^p(X)\), define
\[
  \mathcal O_\alpha(x)
  =
  \frac{1}{p!}\,
  \alpha_{i_1\cdots i_p}(\phi(x))
  \psi^{i_1}(x)\cdots\psi^{i_p}(x).
\]
The Koszul sign rule gives
\[
  Q_A\mathcal O_\alpha=\mathcal O_{d\alpha}.
\]
Therefore local scalar A-model observables are represented by
\[
  H^\bullet_{\rm dR}(X;\mathbb C)
\]
when contact terms are controlled by a specified source and collision
prescription.
The coordinate computation is part of the definition of this local complex:
since \(Q_A\phi^i=\psi^i\) and \(Q_A\psi^i=0\),
\[
  Q_A\mathcal O_\alpha
  =
  \frac{1}{p!}\,
  \partial_j\alpha_{i_1\cdots i_p}(\phi)\,
  \psi^j\psi^{i_1}\cdots\psi^{i_p}.
\]
Antisymmetrizing the coefficient of the product of odd variables gives
\[
  \frac{1}{(p+1)!}(d\alpha)_{j i_1\cdots i_p}
  \psi^j\psi^{i_1}\cdots\psi^{i_p},
\]
which is \(\mathcal O_{d\alpha}\).

\section{The Cauchy--Riemann Section and the Energy Identity}

Let \(J_\Sigma=j_\Sigma\).  For a smooth map \(\phi:\Sigma\to X\), define the
\((0,1)\)-part of its differential by
\[
  \bar\partial_J\phi
  =
  \frac12\bigl(d\phi+J\circ d\phi\circ J_\Sigma\bigr)
  \in
  \Omega^{0,1}(\Sigma,\phi^\ast TX).
\]
The A-model localization equation is
\[
  \bar\partial_J\phi=0.
\]
Its solutions are \(J\)-holomorphic maps.

\paragraph{Pointwise energy identity.}
Fix a conformal metric on \(\Sigma\).  Normalize the norm on the
\((0,1)\)-form bundle by
\[
  |\bar\partial_J\phi|^2
  =
  \frac12
  \bigl|d\phi(e)+Jd\phi(J_\Sigma e)\bigr|^2
\]
for any positively oriented unit vector \(e\), one has the pointwise identity
\[
  \frac12 |d\phi|^2\,d{\rm vol}_\Sigma
  =
  \phi^\ast\omega
  +
  |\bar\partial_J\phi|^2\,d{\rm vol}_\Sigma .
\]
Consequently
\[
  E(\phi)
  :=
  \frac12\int_\Sigma |d\phi|^2\,d{\rm vol}_\Sigma
  =
  \int_\Sigma\phi^\ast\omega
  +
  \int_\Sigma|\bar\partial_J\phi|^2\,d{\rm vol}_\Sigma .
\]

Let \(e_1=e\) and \(e_2=J_\Sigma e\).  Put
\(u=d\phi(e_1)\) and \(v=d\phi(e_2)\).  Compatibility of
\((\omega,J,g_X)\) gives
\[
  \phi^\ast\omega(e_1,e_2)=\omega(u,v)=g_X(Ju,v).
\]
Since \(J\) is \(g_X\)-orthogonal,
\[
  \frac12|u+Jv|^2
  =
  \frac12|u|^2+\frac12|v|^2+g_X(u,Jv)
  =
  \frac12|u|^2+\frac12|v|^2-\omega(u,v).
\]
Multiplying by \(d{\rm vol}_\Sigma\) gives the stated identity.

For a fixed homology class \(\beta\in H_2(X;\mathbb Z)\), the topological
quantity \(\int_\Sigma\phi^\ast\omega=\int_\beta\omega\) is constant on the
component of maps with \(\phi_\ast[\Sigma]=\beta\).  Thus the positive part of
the action is the squared norm of the section \(\bar\partial_J\) plus a
topological term.  In a cohomological representative, the squared norm and its
fermionic linearization arise as \(Q_A\)-variation of a gauge fermion, while
the complexified K\"ahler weight is
\[
  Q^\beta
  :=
  \exp\left(2\pi\ii\int_\beta(B+\ii\omega)\right)
  =
  \exp\left(2\pi\ii\int_\beta B-2\pi\int_\beta\omega\right),
\]
for a closed real \(B\)-field.  This fixes the convention used below for the
Novikov weight \(Q^\beta\).

\section{The Linearized Equation and Virtual Dimension}

Let \(\phi:\Sigma\to X\) solve \(\bar\partial_J\phi=0\).  If \(J\) is
integrable and \(X\) is K\"ahler, the linearization is the Dolbeault operator
\[
  D_\phi
  =
  \bar\partial_{\phi^\ast T^{1,0}X}
  :
  \Gamma(\Sigma,\phi^\ast T^{1,0}X)
  \longrightarrow
  \Omega^{0,1}(\Sigma,\phi^\ast T^{1,0}X),
\]
up to the standard connection expression which is complex-linear on the
holomorphic locus.  For an almost complex target, the linearization is a real
Cauchy--Riemann type Fredholm operator with the same index.

\paragraph{Index of the fixed-domain deformation complex.}
Let \(m=\dim_\mathbb C X\), \(g\) be the genus of \(\Sigma\), and
\(\beta=\phi_\ast[\Sigma]\).  The real Fredholm index calculation for the
deformation
complex
\[
  \Gamma(\Sigma,\phi^\ast TX)
  \xrightarrow{D_\phi}
  \Omega^{0,1}(\Sigma,\phi^\ast TX)
\]
is
\[
  {\rm ind}_{\mathbb R}(D_\phi)
  =
  2\left(m(1-g)+\int_\beta c_1(TX)\right).
\]
For integrable \(J\), the complex index of
\(\bar\partial_{\phi^\ast T^{1,0}X}\) is
\[
  \chi(\Sigma,\phi^\ast T^{1,0}X)
  =
  {\rm rank}_{\mathbb C}(T^{1,0}X)(1-g)
  +
  \deg(\phi^\ast T^{1,0}X)
\]
by Riemann--Roch for holomorphic vector bundles on a compact Riemann surface.
Here the rank is \(m\), and
\[
  \deg(\phi^\ast T^{1,0}X)
  =
  \int_\Sigma \phi^\ast c_1(T^{1,0}X)
  =
  \int_\beta c_1(TX).
\]
The real index is twice the complex index.  For almost complex \(J\), connect
the operator through elliptic real Cauchy--Riemann type operators to a
complex-linear representative; the Fredholm index is locally constant in such
families.

When the complex structure of the domain and \(n\) marked points are also
allowed to vary, the expected real dimension of the stable-map moduli space is
\[
  {\rm vdim}_{\mathbb R}\,\overline{\mathcal M}_{g,n}(X,\beta)
  =
  2\left((1-g)(m-3)+n+\int_\beta c_1(TX)\right).
\]
This is obtained by adding the real dimension \(6g-6+2n\) of the stable
curve moduli space to the fixed-domain index above; the unstable cases are
defined by quotienting by the automorphism group of the marked curve.

\section{Stable Maps and the A-Model Correlator}

\begin{definition}[Stable maps]
Let \(X\) be compact almost K\"ahler.  A genus-\(g\), \(n\)-marked stable map
to \(X\) in class \(\beta\) is a tuple
\[
  (C,p_1,\ldots,p_n,\phi)
\]
where \(C\) is a connected nodal complex curve of arithmetic genus \(g\),
the \(p_a\) are distinct smooth marked points, \(\phi:C\to X\) is continuous
and \(J\)-holomorphic on every irreducible component after normalization, and
\(\phi_\ast[C]=\beta\).  Stability means that every irreducible component on
which \(\phi\) is constant has finite automorphism group after retaining its
marked and nodal special points.
\end{definition}

The compact moduli space is denoted
\[
  \overline{\mathcal M}_{g,n}(X,\beta).
\]
There are evaluation maps
\[
  {\rm ev}_a:\overline{\mathcal M}_{g,n}(X,\beta)\to X,
  \qquad
  {\rm ev}_a(C,p_1,\ldots,p_n,\phi)=\phi(p_a).
\]
The rigorous finite-dimensional object used by the A-model is a virtual
fundamental class
\[
  [\overline{\mathcal M}_{g,n}(X,\beta)]^{\rm vir}
  \in
  H_{{\rm vdim}_{\mathbb R}}\bigl(\overline{\mathcal M}_{g,n}(X,\beta);\mathbb Q\bigr),
\]
or an equivalent virtual integration theory such as a Kuranishi atlas,
polyfold Fredholm structure, or derived geometric perfect obstruction theory
in the algebraic case.  The present chapter records exactly which properties
of this object enter the QFT formulas.

\paragraph{Virtual integration package used below.}
The virtual integration entry \(\mathfrak V_X\) in
Hypothesis~\ref{hyp:a-model-cohomological-sigma-datum} is the following
package.  The symbol
\([\overline{\mathcal M}_{g,n}(X,\beta)]^{\rm vir}\) denotes more than a
homology class attached to a space.  The formulas below require the following
finite-dimensional package.
\[
  \mathfrak V_X
  =
  \bigl(
  [\overline{\mathcal M}_{g,n}(X,\beta)]^{\rm vir},
  {\rm ev}_a,\pi,\rho,\iota_{\Gamma}
  \bigr)_{g,n,\beta,\Gamma}.
\]
Here \({\rm ev}_a\) are the evaluation maps, \(\pi\) denotes the forgetful
maps that drop marked points and stabilize the curve, \(\rho\) denotes
deformation identifications for smooth families of targets or almost complex
structures, and \(\iota_{\Gamma}\) denotes gluing maps associated to nodal
boundary strata labelled by stable graphs \(\Gamma\).  The package is
required to satisfy:
\[
  \dim_{\mathbb R}
  [\overline{\mathcal M}_{g,n}(X,\beta)]^{\rm vir}
  =
  2\left((1-g)(m-3)+n+\int_\beta c_1(TX)\right),
\]
compatibility with pullback of cohomology classes by the evaluation and
forgetful maps, invariance under the deformation identifications \(\rho\),
and compatibility with gluing.  In genus zero, the gluing compatibility used
for associativity is the following boundary-stratum identity.  For a
four-point splitting \(S\) and a decomposition
\(\beta=\beta_1+\beta_2\), let
\[
  \mathcal D_S(\beta_1,\beta_2)
  =
  \overline{\mathcal M}_{0,3}(X,\beta_1)
  \times_X
  \overline{\mathcal M}_{0,3}(X,\beta_2),
\]
where the fiber product uses the two evaluation maps at the node.  The
virtual class on this boundary stratum is required to be
\begin{equation}
  [\mathcal D_S(\beta_1,\beta_2)]^{\rm vir}
  =
  \Delta_X^!\!\left(
  [\overline{\mathcal M}_{0,3}(X,\beta_1)]^{\rm vir}
  \times
  [\overline{\mathcal M}_{0,3}(X,\beta_2)]^{\rm vir}
  \right),
  \label{eq:gw-virtual-splitting-package}
\end{equation}
where \(\Delta_X^!\) is the refined pullback along the diagonal
\(\Delta_X:X\to X\times X\).  If
\(\iota_S:\mathcal D_S(\beta_1,\beta_2)\to
\overline{\mathcal M}_{0,4}(X,\beta)\) is the gluing map, then
\[
  (\iota_S)_*
  \sum_{\beta_1+\beta_2=\beta}
  [\mathcal D_S(\beta_1,\beta_2)]^{\rm vir}
\]
is the virtual class of the boundary fiber over the corresponding boundary
point of
\(\overline{\mathcal M}_{0,4}\cong\mathbb P^1\).  The two boundary points of
\(\overline{\mathcal M}_{0,4}\) used in the WDVV comparison are linearly
equivalent divisors on \(\mathbb P^1\).  This is the nontrivial input from
virtual geometry; it is not derived from the formal algebra of the QFT path
integral in this chapter.

\begin{definition}[Primary Gromov--Witten functional]
Given a virtual integration package \(\mathfrak V_X\) with the properties
specified above, the primary Gromov--Witten functional is the multilinear map
defined for cohomology classes \(\alpha_a\in H^\bullet(X;\mathbb C)\) by
\[
  \left\langle \alpha_1,\ldots,\alpha_n\right\rangle_{g,n,\beta}^{X}
  =
  \int_{[\overline{\mathcal M}_{g,n}(X,\beta)]^{\rm vir}}
  \bigwedge_{a=1}^n {\rm ev}_a^\ast\alpha_a .
\]
The A-model primary correlator at fixed genus is the Novikov series
\[
  \left\langle \alpha_1,\ldots,\alpha_n\right\rangle_{g,n}^{A}
  =
  \sum_{\beta\in H_2(X;\mathbb Z)}
  Q^\beta
  \left\langle \alpha_1,\ldots,\alpha_n\right\rangle_{g,n,\beta}^{X},
\]
whenever the series is defined in the chosen Novikov completion.
\end{definition}

The degree-selection rule is a dimension statement.  The integral vanishes unless
\[
  \sum_{a=1}^n \deg_{\mathbb R}\alpha_a
  =
  {\rm vdim}_{\mathbb R}\,\overline{\mathcal M}_{g,n}(X,\beta).
\]
Equivalently, in complex degree \(|\alpha|_{\mathbb C}=\deg_{\mathbb R}\alpha/2\),
\[
  \sum_{a=1}^n |\alpha_a|_{\mathbb C}
  =
  (1-g)(m-3)+n+\int_\beta c_1(TX).
\]

\section{Quantum Cohomology from Genus-Zero Gluing}

Let \(H=H^\bullet(X;\mathbb C)\), and choose a homogeneous basis
\(\{T_a\}\).  The Poincar\'e pairing is
\[
  \eta_{ab}=\int_X T_a\wedge T_b,
  \qquad
  \eta^{ab}\eta_{bc}=\delta^a_c.
\]

\begin{definition}[Small quantum product]
The genus-zero small quantum product is the Novikov-linear product on
\(H\otimes \Lambda_{\rm Nov}\) determined by
\[
  (T_a\star T_b,T_c)
  =
  \sum_{\beta} Q^\beta
  \left\langle T_a,T_b,T_c\right\rangle_{0,3,\beta}^{X},
\]
where \((\cdot,\cdot)\) is the Poincar\'e pairing.  Equivalently,
\[
  T_a\star T_b
  =
  \sum_{\beta,c,d}
  Q^\beta
  \left\langle T_a,T_b,T_d\right\rangle_{0,3,\beta}^{X}
  \eta^{dc}T_c .
\]
\end{definition}

\paragraph{Associativity from the splitting axiom.}
With the genus-zero splitting identity
\eqref{eq:gw-virtual-splitting-package} as part of the virtual integration
package, the product \(\star\) is associative.  This is a conditional
consequence of the virtual package, not an independent construction of that
package from continuum QFT.

It is enough to prove
\[
  ((T_a\star T_b)\star T_c,T_d)
  =
  (T_a\star (T_b\star T_c),T_d)
\]
for basis elements.  Expanding the left side by the definition of \(\star\)
gives
\[
  \sum_{\beta_1+\beta_2=\beta}
  \sum_{e,f}
  Q^\beta
  \left\langle T_a,T_b,T_e\right\rangle_{0,3,\beta_1}
  \eta^{ef}
  \left\langle T_f,T_c,T_d\right\rangle_{0,3,\beta_2}.
\]
Expanding the right side gives
\[
  \sum_{\beta_1+\beta_2=\beta}
  \sum_{e,f}
  Q^\beta
  \left\langle T_a,T_f,T_d\right\rangle_{0,3,\beta_1}
  \eta^{fe}
  \left\langle T_b,T_c,T_e\right\rangle_{0,3,\beta_2}.
\]
The two expressions are equal because they are the virtual integrals of
\({\rm ev}_1^\ast T_a\,{\rm ev}_2^\ast T_b\,{\rm ev}_3^\ast T_c\,
{\rm ev}_4^\ast T_d\) over two linearly equivalent boundary divisors in
\(\overline{\mathcal M}_{0,4}\cong \mathbb P^1\), pulled back to
\(\overline{\mathcal M}_{0,4}(X,\beta)\).  The splitting axiom identifies the
restriction of the virtual class to the boundary with the fiber product of two
three-point virtual classes over \(X\); the diagonal class of \(X\) is
\[
  [\Delta_X]=\sum_{e,f}\eta^{ef}\,T_e\otimes T_f .
\]
Substituting this diagonal decomposition gives exactly the two displayed
factorizations.  Hence the pairings with all \(T_d\) agree, and
nondegeneracy of \(\eta\) gives associativity.

The calculation shows precisely where the geometry enters: compactness of the
stable-map space, construction of virtual classes, and compatibility of those
classes with the two boundary factorizations of a four-marked rational curve.
The algebraic conclusion is then formal.

\section{Example: Projective Space}

Let \(X=\mathbb P^m\), let \(H\in H^2(\mathbb P^m;\mathbb Z)\) be the
hyperplane class, and normalize
\[
  \int_{\mathbb P^m}H^m=1.
\]
The classical cohomology is \(\mathbb C[H]/(H^{m+1})\).  The small quantum
cohomology is
\[
  QH^\bullet(\mathbb P^m)
  \cong
  \mathbb C[H,Q]/(H^{m+1}-Q),
\]
where \(Q\) is the Novikov variable for the positive line class.

The relation is determined by the three-point invariants.  The expected
complex dimension of
\(\overline{\mathcal M}_{0,3}(\mathbb P^m,d)\) is
\[
  m+(m+1)d.
\]
Thus
\[
  \left\langle H^a,H^b,H^c\right\rangle_{0,3,d}
\]
can be nonzero only when \(a+b+c=m+(m+1)d\).  For \(d=0\), the invariant is
the classical intersection number and gives \(H^a\star H^b=H^{a+b}\) when
\(a+b\le m\).  For \(a+b=m+1\), pair \(H^a\star H^b\) with \(H^{m-c}\).
The only possible correction has \(d=1\) and \(c=0\).  It counts lines in
\(\mathbb P^m\) meeting \(m+1\) independent hyperplane incidence conditions;
equivalently, after choosing a transverse representative, it counts the
unique line determined by two points and the remaining hyperplane conditions.
The result is \(1\).  Hence \(H^{m+1}=Q\).  Associativity then determines all
products.

For \(m=1\), this gives
\[
  QH^\bullet(\mathbb P^1)=\mathbb C[H,Q]/(H^2-Q).
\]
The quantum correction is detected by
\[
  (H\star H,H)
  =
  Q\,\langle H,H,H\rangle_{0,3,1}
  =
  Q,
\]
so \(H\star H=Q\cdot 1\).  This example is the smallest exact check that the
quantum product differs from the classical cup product.

\section{Descendants and the Universal Curve}

Let
\[
  \pi:\mathcal C_{g,n}(X,\beta)\to \overline{\mathcal M}_{g,n}(X,\beta)
\]
be the universal curve and \(s_a\) the section corresponding to the \(a\)-th
marked point.  The cotangent line at the marked point is
\[
  L_a=s_a^\ast\omega_\pi,
\]
where \(\omega_\pi\) is the relative dualizing sheaf.  Its first Chern class
\[
  \psi_a=c_1(L_a)
  \in H^2(\overline{\mathcal M}_{g,n}(X,\beta))
\]
is the gravitational descendant class.  The descendant invariant is
\[
  \left\langle
  \tau_{k_1}(\alpha_1)\cdots \tau_{k_n}(\alpha_n)
  \right\rangle_{g,n,\beta}^{X}
  =
  \int_{[\overline{\mathcal M}_{g,n}(X,\beta)]^{\rm vir}}
  \prod_{a=1}^n
  \psi_a^{k_a}\,{\rm ev}_a^\ast\alpha_a .
\]
In the fixed-source-surface cohomological field theory, descendants arise from
the descent equations of Chapter~\ref{ch:tqft-cohomological-field-theories}.
After coupling to topological gravity, the \(\psi_a\)-classes are the
finite-dimensional remnants of the marked-point cotangent directions.  These
are related constructions occupying different layers: the latter requires
integration over the moduli of domains.

\section{The B-Model Complex}

Let \(X\) be a compact complex \(m\)-fold with \(K_X\simeq\mathcal O_X\), and
choose a holomorphic volume form \(\Omega_X\).  Define the polyvector
Dolbeault complex
\[
  {\rm PV}^{p,q}(X)
  =
  \Omega^{0,q}(X,\wedge^p T^{1,0}X),
  \qquad
  Q_B=\bar\partial .
\]
If
\[
  \mu
  =
  \mu^{i_1\cdots i_p}{}_{\bar j_1\cdots\bar j_q}
  d\bar z^{\bar j_1}\cdots d\bar z^{\bar j_q}
  \otimes
  \partial_{i_1}\wedge\cdots\wedge\partial_{i_p},
\]
then the corresponding local observable is represented by
\[
  \mathcal O_\mu
  =
  \mu^{i_1\cdots i_p}{}_{\bar j_1\cdots\bar j_q}(\phi)
  \eta^{\bar j_1}\cdots\eta^{\bar j_q}
  \theta_{i_1}\cdots\theta_{i_p},
\]
where \(\eta^{\bar j}\) is the odd variable corresponding to
\(d\bar z^{\bar j}\) and \(\theta_i\) is the odd variable dual to
\(\partial_i\).  The scalar \(Q_B\)-cohomology of local observables is
\[
  H^q(X,\wedge^pT^{1,0}X).
\]

The holomorphic volume form converts polyvectors into forms by contraction,
\[
  \iota_\mu\Omega_X\in \Omega^{m-p,q}(X).
\]
For \(\mu_1,\ldots,\mu_k\) with total polyvector degree \(m\) and total
Dolbeault degree \(m\), the genus-zero classical B-model pairing is
\[
  \left\langle \mu_1,\ldots,\mu_k\right\rangle_B
  =
  \int_X
  \Omega_X\wedge
  \iota_{\mu_1\wedge\cdots\wedge\mu_k}\Omega_X .
\]
If either total degree condition fails, the integral is zero by form degree.

\begin{remark}[B-model dependence on complex structure]
Infinitesimal deformations of the complex structure are represented by
\[
  \mu\in \Omega^{0,1}(X,T^{1,0}X).
\]
The integrability equation for a finite deformation is the
Maurer--Cartan equation
\[
  \bar\partial\mu+\frac12[\mu,\mu]=0,
\]
where \([\cdot,\cdot]\) is the Schouten--Nijenhuis bracket extended to
Dolbeault forms.  Thus the B-model deformation problem is controlled by the
differential graded Lie algebra
\[
  \bigl({\rm PV}^{1,\bullet}(X),\bar\partial,[\cdot,\cdot]\bigr).
\]

Write the deformed \((0,1)\)-operator on functions as
\[
  \bar\partial_\mu=\bar\partial+\mu\cdot\partial .
\]
A complex structure is integrable exactly when
\(\bar\partial_\mu^2=0\).  Expanding the square on local functions gives
\[
  \bar\partial\mu+\frac12[\mu,\mu],
\]
with the bracket equal to the Schouten--Nijenhuis bracket in the vector index
and the wedge product in the Dolbeault indices.  Hence integrability is the
displayed Maurer--Cartan equation.
\end{remark}

The B-twist uses the axial \(U(1)\) R-symmetry of the parent
\(\mathcal N=(2,2)\) sigma model.  The anomaly of this symmetry is measured
by \(c_1(TX)\), so the closed B-model requires \(c_1(TX)=0\), equivalently
triviality of \(K_X\) after choosing the holomorphic volume form needed for
the trace pairing.  Boundary conditions and open B-models require additional
Chan--Paton and anomaly data.

\section{A-Model and B-Model Parameter Dependence}

The two models organize target geometry in complementary ways.

\begin{itemize}
\item The A-model local scalar complex is \(H^\bullet_{\rm dR}(X)\).  Its
  instanton weights depend on the complexified K\"ahler class
  \([B+\ii\omega]\) through \(Q^\beta\).  A change of the target complex
  structure changes the representative of the Cauchy--Riemann section, but
  under a \(Q_A\)-preserving regulator and compactness hypotheses the induced
  variation of primary correlators is \(Q_A\)-exact.
\item The B-model local scalar complex is
  \(H^q(X,\wedge^pT^{1,0}X)\).  Its deformation theory is controlled by the
  complex-structure Maurer--Cartan equation above.  K\"ahler variations occur
  in \(Q_B\)-exact metric-response terms, again provided the regulator and
  integration cycle preserve the \(Q_B\)-Ward identity.
\end{itemize}

These statements are Ward-identity statements.  They require the hypotheses
listed in Chapter~\ref{ch:tqft-cohomological-field-theories}: a
\(Q\)-invariant regulated cycle, control of boundary terms at noncompact
ends, and a treatment of contact terms for colliding insertions.

\section{Continuum QFT Status}

The A-model has a mathematically rich finite-dimensional output in
Gromov--Witten theory.  The B-model has a finite-dimensional classical closed
sector expressed through polyvector Dolbeault cohomology and, at higher
genus, through the Kodaira--Spencer/BCOV effective theory.  These
outputs are indispensable, and the continuum sigma-model QFT still demands a
field-theoretic construction.

\begin{openproblem}[Topological sigma model as continuum QFT]
\label{op:topological-sigma-model-continuum-qft}
Construct A- and B-models as continuum cohomological QFTs with explicit
regulated field spaces, \(Q\)-invariant integration cycles, anomaly lines,
compactified zero-mode spaces, compatible virtual integration, and gluing
maps for cutting and sewing source surfaces.  The construction should recover the
Gromov--Witten and B-model finite-dimensional formulas as theorems rather
than as definitions of the path integral.
\end{openproblem}

This open problem is the precise interface where enumerative geometry,
supersymmetric sigma-model regularization, and functorial QFT meet.
